% Sections
\newcommand{\langexperience}{EXPERIENCIA}
\newcommand{\langeducation}{EDUCACIÓN}
\newcommand{\langcertificates}{CERTIFICADOS}
\newcommand{\langprojects}{PROYECTOS}
\newcommand{\langskills}{HABILIDADES}
\newcommand{\langgreetings}{SALUDOS}
\newcommand{\langmetadata}{
	Currículum creado para reflejar únicamente mi experiencia como Desarrollador en la Nube.

	Para un resumen completo de mis habilidades, ¡Consulta mis otros currículums específicos para cada tecnología!
}

% Contact
\newcommand{\langpersonA}{Jesús Antonio González Quevedo}
\newcommand{\langpersonB}{https://maps.app.goo.gl/U12Lxq6ALWDSuLdV9}
\newcommand{\langpersonC}{Puebla, Puebla}
\newcommand{\langpersonD}{tel:522222650774}
\newcommand{\langpersonE}{+52 222 265 0774}
\newcommand{\langpersonF}{mailto:jesugq@yahoo.com}
\newcommand{\langpersonG}{jesugq@yahoo.com}
\newcommand{\langpersonH}{https://www.linkedin.com/in/jesugq}
\newcommand{\langpersonI}{jesugq}

% Cover
\newcommand{\langcoverA}{¡Saludos!}
\newcommand{\langcoverB}{¡Adiós!}

% Experience
\newcommand{\langmotivusA}{Motivus}
\newcommand{\langmotivusB}{https://motivus.com/}
\newcommand{\langmotivusC}{faGlobe}
\newcommand{\langmotivusD}{Ago 2021 — Abr 2025}
\newcommand{\langmotivusE}{Desarrollador Asociado}
\newcommand{\langmotivusF}{Google Cloud Platform, Microservicios, Bases de Datos, Git}
\newcommand{\langmotivusG}{Trabajo Remoto}
\newcommand{\langmotivusH}{
	\item Lideré un equipo de 10 como Scrum master y Desarrollador para una plataforma de evaluación de reclutas.
	\item Agregué la habilidad de buscar las universidades, grupos, exámenes y calificaciones de los reclutas en la base de datos.
	\item Creé las vistas de usuario, usuarios por grupo, gráficas de calificación dinámicas por grupo, búsqueda, etc.
	\item Supervisé los cambios en el repositorio de control de versiones en GitHub para nuestro equipo de backend y frontend.
	\item Participé en un rebase de git para arreglar los conflictos entre las ramas de desarrollo de múltiples equipos.
	\item Fui el mayor contribuyente al repositorio Git de Java, dando el 60\% de de líneas de código nuevas totales.
	\item Recibí una licencia de Udemy (en uso!) por años consecutivos de alto rendimiento (Evaluación: 4 — Excede Expectativas.)
	\item Me uní a un equipo de incidentes de producción en Citibanamex  a través de la firma de consultoría Motivus.
}
%
\newcommand{\langcitibanamexA}{Banca Empresarial Citibanamex (Contratante)}
\newcommand{\langcitibanamexB}{https://www.banamex.com}
\newcommand{\langcitibanamexC}{faGlobe}
\newcommand{\langcitibanamexD}{Dic 2021 — Feb 2025}
\newcommand{\langcitibanamexE}{Desarrollador QA Asociado}
\newcommand{\langcitibanamexF}{Google Cloud Platform, Microservicios, Bases de Datos, Git}
\newcommand{\langcitibanamexG}{Trabajo Remoto}
\newcommand{\langcitibanamexH}{
	\item Supervisé la liberación de 231 actualizaciones de microservicios de 32 equipos de desarrollo durante un año.
	\item Aprendí la funcionalidad de 91 diferentes microservicios a través 5 ambientes y 6 gateways.
	\item Verifiqué la comunicación JDBC entre microservicios y sus bases de datos SQL/NoSQL, 61 scripts de consulta.
		\item Analicé 6 incidentes de producción rastreando peticiones a través de microservicios gracias a sus registros encriptados.
	\item Ocasionalmente programé en pareja para encontrar las razones detrás de las discrepancias en formato de errores.
	\item Ocasionalmente navegamos a través de commits de Git para encontrar la línea de código que reintegró un bug ya arreglado.
	\item Ocasionalmente me conecté a través de SSH a un servidor Linux para actualizar archivos de configuración y correr bash scripts.
}

% Education
\newcommand{\langtecnologicomonterreyA}{Instituto Tecnológico de Estudios Superiores de Monterrey (ITESM PUE)}
\newcommand{\langtecnologicomonterreyB}{https://tec.mx/}
\newcommand{\langtecnologicomonterreyC}{faGlobe}
\newcommand{\langtecnologicomonterreyD}{https://certificados.tec.mx/certificate/dd14d13b16e750099a9783188f343e25}
\newcommand{\langtecnologicomonterreyE}{faCheckCircle}
\newcommand{\langtecnologicomonterreyF}{Cédula Profesional de Ingeniería en Tecnologías Computacionales}
\newcommand{\langtecnologicomonterreyG}{Identity Platform, App Engine, Linux, Git}
\newcommand{\langtecnologicomonterreyH}{
	\item Empecé a utilizar Linux en 2018 y lo he estado usando desde entonces (Ubuntu, Mint, Kali, Arch, Manjaro, Fedora).
	\item Utilicé Google Cloud Platform crear un de inicio de sesión con el Identity Platform para el curso de Lab. Desarrollo de Aplicaciones.
	\item Utilicé Google Cloud Platform para crear una plataforma de exámenes en línea para el curso de Desarrollo de Aplicaciones Web.
	\item Utilicé Git en la mayoría de los proyectos de programación (no listados por el uso de tecnologías externas a la Nube).
}

% Certificates
\newcommand{\bodyclouddeveloperA}{Google Cloud Platform}
\newcommand{\bodyclouddeveloperB}{https://cloud.google.com/learn/certification/cloud-developer}
\newcommand{\bodyclouddeveloperC}{faGlobe}
\newcommand{\bodyclouddeveloperD}{https://www.credential.net/f05088e4-157e-40be-8d4f-b918be06f6c9}
\newcommand{\bodyclouddeveloperE}{faCheckCircle}
\newcommand{\bodyclouddeveloperF}{Desarrollador Profesional en la Nube (384/384 lecciones)}
\newcommand{\bodyclouddeveloperG}{Docker, Compute Engine, Cloud Datastore, Pub/Sub, Kubernetes Engine}
\newcommand{\bodyclouddeveloperH}{
	\item Desplegué un ambiente de desarrollo Java en una máquina virtual Compute Engine con librerías de software preinstaladas.
	\item Desplegué una base de datos Cloud Datastore NoSQL provisionado por App Engine para guardar respuestas de un examen en línea.
	\item Desplegué autenticación de usuario con Identity Platform para restringir el accesso a Cloud Storage y Cloud Datastore.
	\item Desplegué un servicio de backend Cloud Spanner y Cloud Pub/Sub conectado a Natural Language API para analizar los datos.
	\item Desplegué y configuré un proyecto de Kubernetes Engine utilizando una imagen de Docker de muestra con Cloud Build.
	\item Configuré el despliegue de Kubernetes Engine para otorgar escalabilidad además de despliegues Azul-Verde y Canario.
	\item Vea la lista completa de insignias (la primera mitad hace referencia a este curso) aquí. \hypericon{https://www.cloudskillsboost.google/public_profiles/f0976fc9-283a-49aa-b9a1-a214186ac3e4}{faExternalLink}
}
%
\newcommand{\bodycloudengineerA}{Google Cloud Platform}
\newcommand{\bodycloudengineerB}{https://cloud.google.com/learn/certification/cloud-engineer}
\newcommand{\bodycloudengineerC}{faGlobe}
\newcommand{\bodycloudengineerD}{https://www.credential.net/59c21655-2c0d-4803-ba06-3bad73fa3b69}
\newcommand{\bodycloudengineerE}{faCheckCircle}
\newcommand{\bodycloudengineerF}{Ingeniero Asociado en la Nube (306/306 lecciones)}
\newcommand{\bodycloudengineerG}{Docker, App Engine, Cloud Storage, BigQuery, Kubernetes Engine, Cloud Run}
\newcommand{\bodycloudengineerH}{
	\item Exploré las características de Google Cloud Console; Mercado, Redes, Máquinas Virtuales, Manejo de Cargas, Monitoreo.
	\item Desplegué una solución Cloud Storage con listas de acceso, claves CSEK rotativas, versiones de archivos y gestión del ciclo de vida.
	\item Importé registros de servidor de Cloud Storage a BigQuery y realicé consultas a la información desde la Cloud Console.
	\item Desplegué la misma imagen de Docker de muestra en App Engine, Kubernetes Engine, y Cloud Run.
	\item Desplegué una pipeline básica de DevOps usando Cloud Build para automatizar despliegues tras cada commit a Source Repositories.
	\item Vea la lista completa de insignias (la segunda mitad hace referencia a este curso) aquí. \hypericon{https://www.cloudskillsboost.google/public_profiles/f0976fc9-283a-49aa-b9a1-a214186ac3e4}{faExternalLink}
}

% Skills
\newcommand{\langlanguagesA}{Lenguajes}
\newcommand{\langlanguagesB}{Inglés Avanzado, Español Nativo}
%
\newcommand{\langdevelopmentA}{Development}
\newcommand{\langdevelopmentB}{Google Cloud Platform, Cloud Console, Cloud Shell, Git}
%
\newcommand{\langdevelopmentC}{}
\newcommand{\langdevelopmentD}{Linux (Fedora Linux), Github, Bitbucket, Source Repositories, TectiaSSH}
%
\newcommand{\langconceptsA}{Conceptos}
\newcommand{\langconceptsB}{Desarrollo en la Nube, Prog. Orientada a Objetos, Patrones de Diseño, Ctrl. de Ver., Microservicios, Bases de Datos}
%
\newcommand{\langconceptsC}{}
\newcommand{\langconceptsD}{Gang of Four, SOLID, ACID, DDD, MVC, ORM, IoC, DI, CI/CD, JVM, }
%
\newcommand{\langprinciplesA}{Principles}
\newcommand{\langprinciplesB}{Autenticación, Autorización, Despliegues Azul-Verde, Contenerización, Despliegues Canario}
%
\newcommand{\langprinciplesC}{}
\newcommand{\langprinciplesD}{Funciones, Manejo de Identidad, Imágenes, Balanceo de Carga, Publicador / Subscriptor, Máquinas Virtuales}